\section{A.6: stable torsion and a characteristic-zero lift}
\label{sec:A-lifting}

\begin{theorem}[The special-fibre identity and the specialization obstruction]\label{thm:A-special-torsion}
\status{Known, with the derivation below.} For every smooth projective
$Z/k$, with $k$ perfect, there is a natural Galois-equivariant identity
after extension to $\bar k$:
\begin{equation}\label{eq:A-special-torsion}
 T^j(Z_{\bar k})[p]_{\ss}\simeq
 H^j_{\et}(Z_{\bar k},\Zp)[p]\otimes_{\Fp}\bar k.
\end{equation}
Let $\mathcal Z/\OO_K$ be a smooth proper lift, where $\OO_K$ is a
complete mixed-characteristic discrete valuation ring with perfect
residue field $k$ and absolute ramification index $e$.  Specialization
is a map of continuous Galois modules
\[
 s^j:H^j_{\et}(Z_{\bar k},\Zp)\longrightarrow
 H^j_{\et}(\mathcal Z_{\bar K},\Zp).
\]
It need not be injective on torsion.  More precisely, put
$\mathfrak S=W(k)[[u]]$, with $u$ mapping to a chosen uniformizer, and
\[
 M^j=H^j_\Delta(\widehat{\mathcal Z}/\mathfrak S)[u^\infty].
\]
Then the Breuil--Kisin prismatic Frobenius gives
\begin{equation}\label{eq:A-prismatic-kernel}
 \ker s^j\simeq
 \left((M^j/uM^j)\otimes_{W(k)}W(\bar k)\right)^{\varphi=1}.
\end{equation}
In particular $s^j$ is injective if $e(j-1)<p-1$; its kernel is killed
by $p^{j-1}$ if $e(j-1)<2(p-1)$, and by $p$ at the boundary
$e(j-1)=p-1$.
\end{theorem}
\begin{proof}
Apply the natural integral identification in
\citet[Corollary 4.11]{Ogus1975} and take the kernel of multiplication
by $p$.  Exactness of stable extraction and flatness of $W(\bar k)$
give \eqref{eq:A-special-torsion}.  This argument includes torsion;
rational slope comparison alone would not suffice.

For the additional lift, the source and target of specialization are
the $\varphi=1$ complexes before and after inverting the prismatic
parameter, respectively.  The comparison is
\citet[Theorem 4.10]{LiLiu2022}.  Their Lemma 4.11 makes
$\varphi-1$ surjective on cohomology.  Hence taking its kernel in the
localization diagram identifies $\ker s^j$ with the Frobenius-fixed
parameter-torsion module.  The base change and reduction in their
Theorem 4.14 identify it with the right side of
\eqref{eq:A-prismatic-kernel}; their Corollary 4.15 gives the stated
bounds.  These theorems apply to the formal completion of our smooth
proper lift, with the same ramification index and coefficient lattice.
Thus we have specified both the exact special-fibre torsion group and
the obstruction to transporting it injectively to the generic fibre.
\end{proof}

\begin{theorem}[Refutation of the unrestricted lifted-torsion criterion]\label{thm:A-BMS-counterexample}
\status{Refuted, by a complete construction.} There is a smooth
projective surface $X_0/\mathbf F_2$, with a polarized smooth
projective lift $\mathcal X/\mathbf Z_2$, such that
\[
 H^2_{\et}(\mathcal X_{\overline{\mathbf Q}_2},\mathbf Z_2)
 \text{ is torsion-free},\qquad
 T^2(X_{0,\overline{\mathbf F}_2})[2]_{\ss}\simeq
 \overline{\mathbf F}_2\ne0.
\]
Consequently the vanishing in 8.5.10 cannot in general be replaced by
torsion-freeness of the analogous generic-fibre cohomology group.
This is the phenomenon of \citet[Theorem 2.1 and Remark 2.3]{BMS2016};
we give an explicit choice of the cover and of its lift in their
construction.  We claim no priority.
\end{theorem}
\begin{proof}
\emph{The cover and its involution.}
In $\mathbf P^5_{\mathbf Z_2}$ with coordinates
$(x_1,x_2,x_3,y_1,y_2,y_3)$, let $Q$ be cut out by the three quadrics
\begin{equation}\label{eq:A-K3-quadrics}
 x_i^2+x_jx_k-y_i^2-x_iy_i=0\quad
 ((i,j,k)=(1,2,3),(2,3,1),(3,1,2)).
\end{equation}
The involution
$\iota(x_i,y_i)=(x_i,-x_i-y_i)$ preserves each equation, by direct
substitution.  Modulo two the Jacobian has rank three everywhere on
the complete intersection.  Here is an elementary verification.
At a hypothetical singular point, the minor in the three $y$ columns
forces $x_1x_2x_3=0$.  By cyclic symmetry suppose $x_1=0$.
The minor in columns $x_1,x_2,y_2$ gives $x_2^2x_3=0$.
If $x_2=0$, the first two equations give $y_1=y_2=0$, and the minor
in the three $x$ columns gives $x_3^2y_3=0$.  The last equation then
gives $x_3=y_3=0$, impossible in projective space.  The case $x_3=0$
is the same after interchanging the remaining indices.
Thus the special fibre is a smooth complete intersection surface.
The equations form a relative regular sequence and give a flat
projective lift.  Its nonsmooth locus is closed and proper over
$\mathbf Z_2$ and misses the closed fibre, so is empty.

In characteristic two a projective fixed point of $\iota$ has
$x_1=x_2=x_3=0$, and the equations force all the $y_i$ to vanish.
Thus the action is free on the closed fibre.  Its fixed subscheme is
proper, so it is empty over $\mathbf Z_2$ as well.  The quotient
$S=Q/\langle\iota\rangle$ therefore exists as a smooth projective
scheme and $Q\to S$ is a finite etale double cover.  For projectivity,
tensor an ample line bundle with its translate, use the resulting
group linearization, and descend an appropriate positive power.
The geometric fibres of $Q$ are connected complete intersections
with trivial canonical bundle.  They have trivial etale fundamental
group: apply the projective Lefschetz theorem successively to the
complete intersection flag, using the local complete intersection
purity condition in \citet[Expos\'e XII, Corollary 3.5]{SGA2}.
If an intermediate member of the displayed flag is singular, replace
the flag by general linear combinations of the three equations;
the final smooth complete intersection is unchanged, and the same
local complete intersection version applies.  Consequently
\[
 \pi_1(S_{\overline{\mathbf F}_2})=\mathbf Z/2\mathbf Z.
\]

\emph{A translation that becomes zero.}
Take the smooth elliptic scheme
\[
 E:\quad y^2+xy=x^3+1
\]
over $\mathbf Z_2$.  Its discriminant is $-433$, a unit.  Its
characteristic-two special fibre is ordinary, since the Hasse
invariant is $a_1=1$.  There is a unique odd $u\in\mathbf Z_2$ with
$u^3+u^2+64=0$: the polynomial vanishes modulo two at $u=1$ and its
derivative is odd there, so Hensel's lemma applies.  Put
\[
 P=(u/4,-u/8)\in E(\mathbf Q_2).
\]
The defining equation reduces to precisely the polynomial for $u$,
and the elliptic inverse of $(x,y)$ is $(x,-y-x)$, so $P=-P$ and
$P\ne O$.  Properness extends $P$ to a section over $\mathbf Z_2$.
In projective coordinates it is $[2u:-u:8]$, whose reduction is
$[0:1:0]=O$.  Hence the homomorphism
$\eta:(\mathbf Z/2)_{\mathbf Z_2}\to E$, $1\mapsto P$, is nonzero
generically and zero on the special fibre.

\emph{The threefold.}
Let
\[
 D=(Q\times E)/\langle(\iota,+P)\rangle.
\]
The action is free because the action on $Q$ is free.  Thus $D$ is
smooth projective over $\mathbf Z_2$, of relative dimension three.
Its special fibre is $S_0\times E_0$, because $P_0=O$.
Its characteristic-zero geometric fibre maps to
$E'=E/\langle P\rangle$ with simply connected fibre $Q$.
The homotopy exact sequence for this smooth proper fibration gives
\[
 \pi_1(D_{\overline{\mathbf Q}_2})\simeq\pi_1(E')\simeq
 \widehat{\mathbf Z}^{\,2}.
\]
Equivalently, this follows by comparing the two double-cover exact
sequences for $Q\times E\to D$ and $E\to E'$.
Its maximal abelian pro-two quotient is $\mathbf Z_2^2$.
The integral universal coefficient theorem, or singular comparison
after embedding a field of definition into $\mathbf C$, now gives
torsion-freeness of $H^2_{\et}(D_{\overline{\mathbf Q}_2},\mathbf Z_2)$.

On the other hand,
$H^1_{\et}(S_{0,\bar k},\mathbf Z_2)=0$ and
$H^1_{\et}(S_{0,\bar k},\mathbf F_2)=\mathbf F_2$.
The coefficient Bockstein therefore produces a nonzero class of
order two in $H^2_{\et}(S_{0,\bar k},\mathbf Z_2)$.
For $D_0=S_0\times E_0$, the decomposition of its fundamental group
gives
\[
 H^1_{\et}(D_{0,\bar k},\mathbf Z_2)\simeq\mathbf Z_2,
 \qquad H^1_{\et}(D_{0,\bar k},\mathbf F_2)\simeq\mathbf F_2^2;
\]
the ordinary elliptic curve contributes the free summand.  The
coefficient sequence thus gives
$H^2_{\et}(D_{0,\bar k},\mathbf Z_2)[2]\simeq\mathbf F_2$.

\emph{Selection of the surface, with no unspecified existence step.}
Choose a projective embedding of $D$ over $\mathbf Z_2$ by the
descended ample bundle above and a sufficiently large power.  The
construction can be made deterministic: order positive powers,
choose the first that is very ample, choose a basis in a fixed
monomial order, and enumerate homogeneous polynomials over
$\mathbf F_2$ by degree and lexicographic coefficient order.
Start at degree $d_{\rm int}(D_0)$ from
Theorem~\ref{thm:A-high-degree}, and choose the first polynomial
whose zero locus in $D_0$ is smooth.  Such a polynomial exists by
\citet[Theorem 1.1]{Poonen2004}; smoothness is decidable by the
Jacobian criterion and saturation on finitely many affine charts.
Lift its coefficients to $0,1\in\mathbf Z_2$ and call the resulting
divisor $\mathcal X\subset D$.  Its special fibre is smooth and its
defining section is regular.  The same properness and flatness
argument as above proves that $\mathcal X$ is smooth projective over
$\mathbf Z_2$; its polarization is the restricted embedding.

Integral weak Lefschetz in characteristic zero makes the map on
$H^2_{\et}(\mathbf Z_2)$ injective with torsion-free cokernel, so
$H^2_{\et}(\mathcal X_{\overline{\mathbf Q}_2},\mathbf Z_2)$ is
torsion-free.  In characteristic two,
Lemma~\ref{lem:A-R1} and the integral argument in
Theorem~\ref{thm:A-unit-count} give isomorphisms on $H^1_{\et}$ with
both $\mathbf Z_2$ and $\mathbf F_2$ coefficients.  Naturality of the
Bockstein therefore gives
\[
H^2_{\et}(X_{0,\bar k},\mathbf Z_2)[2]\simeq\mathbf F_2.
\]
Equation~\eqref{eq:A-special-torsion} proves the asserted nonzero
stable crystalline torsion.  All hypotheses in A.6, including a
lifted polarization and dimension two, have now been verified.
\end{proof}

\begin{remark}[Scope of the computer verification]\label{rem:A-K3-computation}
\status{Proved}
The program \path{code/k3_lift_check.py} verifies the integral
involution identities, computes the unit Groebner ideal of the
singular locus in every one of the six projective charts over
$\mathbf F_2$, checks $\Delta_E=-433$, and computes the Hensel root
through precision $2^{32}$.  Its output is
\path{outputs/k3_lift_check.json}.  These are exact polynomial
computations.  We do not claim to have executed the final surface
enumeration: its termination and the cohomological assertions are
proved above.  The construction is a specified terminating
construction, not a printed equation for that last ample section.
\end{remark}

\begin{corollary}[A valid restricted implication]\label{cor:A-lifting-small-range}
\status{Proved under stated hypotheses.} If the lift in
Theorem~\ref{thm:A-special-torsion} satisfies $e(n-1)<p-1$ and its
generic $H^n_{\et}(\Zp)$ is torsion-free, then
$T^n(Z_{\bar k})[p]_{\ss}=0$.  For a lift over $W(k)$ this condition
is $n<p$.  The example above has $p=n=2$ and lies exactly on the
boundary.  No converse is asserted by this corollary.
\end{corollary}
\begin{proof}
In the indicated range specialization is injective.  Its source
cannot contain $p$-torsion if its target contains none.  Apply
\eqref{eq:A-special-torsion}.
\end{proof}
